\section{Introduction}
A provider that serves heterogeneous histories may have to commit to a small common menu of terminal commands. The provider can adjust service before sending a command, but the terminal interface cannot recover the history from an uncharged side channel. Choosing the menu and allocating the accepted obligation are therefore coupled decisions. A command that improves terminal reward can require an opening payment, become unsafe in some histories, or displace another command from the interface. Optimizing the allocation after fixing the menu does not resolve this joint design problem.

We study finite-catalog renewal design with individual expected-obligation caps, history-specific realization ceilings, convex intermediate-service costs, and additive command-opening charges. The provider must meet the accepted obligation exactly and retain the original command budget. We establish a selected-boundary decomposition that turns the joint problem into an ordered resource-constrained path problem. A path chooses the commands. Each adjacent selected pair owns a terminal-reward increment and the service costs of precisely the histories whose eligibility ends at that pair. A terminal edge accounts for the remaining histories. Unselected catalog entries do not divide these ownership sets.

The decomposition yields an additive approximation scheme whose polynomial degree does not depend on the number of eligibility classes. A scalar resource grid replaces the earlier search over class means. Rounding an optimal path loses a controlled amount of resource; the selected path always has enough capacity to restore the exact obligation. This capacity identity is essential: a resource-grid solution alone would only solve a perturbed-obligation problem. We recover a feasible policy at the original obligation and certify its distance from the globally optimal finite-catalog value. The guarantee is additive, not relative, and does not enlarge the installed menu.

The result complements, rather than replaces, exact common-eligibility pooling. Histories that admit the same catalog prefix retain their individual caps and service costs under pooling. Their correction depends only on the last eligible selected command. With one eligibility class and rational quadratic primitives, a polynomial ordered-path calculation solves the joint problem exactly, including interior obligations and zero service curvatures. With arbitrary eligibility, the new resource path provides polynomial dependence on inverse normalized accuracy instead of a class-dependent exponent. We separate these two computational statements throughout.

Catalog dependence also becomes explicit. Inserting an unused command can split a full-catalog eligibility class without changing the optimal decision. By contrast, the resource contribution of an edge between two retained selected commands is unchanged by such insertion. We prove insertion and threshold-margin results, identify a discontinuity at a threshold crossing, and give robust value brackets for uncertain ceilings. These results distinguish genuine changes in the feasible interface from changes in a convenient partition.

A uniform deletion envelope further separates harmful refinement from irrelevant candidate proliferation. It bounds the payoff contribution of a non-anchor command in every possible selected neighboring context by a single minimum-command expression. A command whose opening charge covers this envelope can be removed without changing any individual target or the optimal value. All certified removals can be made together, and the bounds do not depend on other non-anchor candidates. A checked global interval on the reduced catalog then transfers to the original catalog through a separately verified exclusion record. This is an optimization guarantee, not a conclusion drawn from observing an unchanged incumbent.

Our analysis uses classical tools with specific roles. Separable convex allocation supplies service values and marginal-price witnesses; \citet{Patriksson2008} surveys this problem and its algorithms. Resource discretization has a substantial history in constrained shortest-path approximation, including \citet{Hassin1992} and \citet{LorenzRaz2001}. Totally monotone matrix search provides efficient convolution and independent Bellman verification \citep{Aggarwal1987}. We do not present these components as new. The contribution is the exact selected-boundary reduction, the capacity identity that permits original-obligation recovery, and their consequence for finite-catalog joint design. The maximum over installed menus need not be concave, even though every fixed-menu response and every arc kernel is concave.

Finite policy graphs approximate continuation contracts by a finite recursive representation \citep{Zhang2012}. Our restriction instead concerns the shared terminal alphabet; the pre-draw allocation and the original equality remain explicit. Adjacent lotteries are related to mean-preserving random splitting in dual quantization \citep{PagesWilbertz2012}, while ordered quantization can exploit matrix searching \citep{Wu1991}. Those connections motivate the representation, but do not by themselves resolve opening charges, history-specific safety prefixes, and endogenous allocation of the accepted obligation. The selected-boundary reduction addresses that joint decision.

The computational study is designed around the structural questions. It includes heterogeneous ceilings within strict interior prefixes, controlled eligibility splits, catalog-refinement paths, provably unused splitting commands, and every previously reported difficult instance. Exact enumeration supplies small-instance references. New resource-grid certificates are checked through a separate marginal-price construction and a different matrix-search algorithm. Matched-machine group and individual subdivisions, price-set sensitivity, and numerical mixed-integer envelopes expose the costs as well as the benefits of certification. The study reports certificate sizes, rational bit lengths, relative gaps, and unsuccessful historical requests. All inputs are synthetic; no claim of field calibration or general solver dominance is made.

The same capacity-repair argument applies to compatible production-module configurations with a divisible order and concave operating returns. We state this second model independently, without renewal histories or command lotteries. The archived continuous-design, institutional, and coarsening results remain available with their complete proofs and original assumptions. The present article concentrates on the finite-catalog joint design problem and its strengthened guarantees.
